\documentclass[11pt,a4paper]{article}
\usepackage[margin=25mm,headheight=15pt]{geometry}
\usepackage[T1]{fontenc}
\usepackage[utf8]{inputenc}
\usepackage{lmodern,microtype}
\usepackage{amsmath,amssymb,amsthm,mathtools}
\usepackage{booktabs,longtable,tabularx,array}
\usepackage{xcolor,fancyhdr}
\usepackage[round,authoryear]{natbib}
\usepackage[breaklinks,colorlinks=true,linkcolor=black,citecolor=black,urlcolor=blue!50!black]{hyperref}
\usepackage{url}
\urlstyle{same}
\setlength{\parindent}{0pt}
\setlength{\parskip}{5pt}
\setlength{\emergencystretch}{3em}
\renewcommand{\arraystretch}{1.18}
\pagestyle{fancy}\fancyhf{}
\fancyhead[L]{\small Liquidity-aware hedging / Research working paper}
\fancyhead[R]{\small AMS 518 \& AMS 520}
\fancyfoot[L]{\footnotesize Theory, literature, and reproducible design}
\fancyfoot[R]{\thepage}
\newcommand{\E}{\mathbb E}
\newcommand{\Pp}{\mathbb P}
\newcommand{\Q}{\mathbb Q}
\newcommand{\F}{\mathcal F}
\newcommand{\G}{\mathcal G}
\newcommand{\ind}{\mathbf 1}
\newcommand{\CVaR}{\operatorname{CVaR}}
\newcommand{\VaR}{\operatorname{VaR}}
\newcommand{\Var}{\operatorname{Var}}
\newcommand{\Cov}{\operatorname{Cov}}
\newcommand{\argmin}{\operatorname*{arg\,min}}
\newcommand{\norm}[1]{\left\lVert#1\right\rVert}
\newtheorem{proposition}{Proposition}[section]
\newtheorem{remark}[proposition]{Remark}
\newcolumntype{Y}{>{\raggedright\arraybackslash}X}
\hypersetup{pdftitle={Liquidity-Aware Derivative Hedging: Theoretical Analysis and Literature Review},pdfauthor={Prepared for Michael Perez; AI-assisted research draft},pdfsubject={AMS 518 and AMS 520 project research foundation}}
\begin{document}
\begin{titlepage}
\thispagestyle{empty}
{\small RESEARCH WORKING PAPER \hfill 23 September 2026}\par
\vspace{22mm}
{\Huge\bfseries Liquidity-Aware\\[3mm] Derivative Hedging}\par
\vspace{8mm}
{\Large Order-Book Learning, CVaR Optimization,\\[2mm] and Local--Stochastic Volatility}\par
\vspace{9mm}
{\large Theoretical analysis and critical literature review}\par
\vspace{15mm}
\par\medskip
\begin{tabularx}{\textwidth}{@{}p{34mm}Y@{}}
Prepared for & Michael Perez\\
Academic context & AMS 518: Advanced Stochastic Models, Risk Assessment, and Portfolio Optimization; AMS 520: Machine Learning in Quantitative Finance\\
Research platform & \texttt{michaeletro/Derivative\_Trading\_Strategy}\\
Implementation & Receipt-time liquidity targets, chronological validation, and persistence/residual-ridge baselines\\
Status & AI-assisted project draft for review; not a completed empirical study or an instructor-approved combined submission
\end{tabularx}
\par\medskip
\vfill
\hrule\vspace{4mm}
\textbf{Evidence boundary.} The accompanying notebook executes on explicitly synthetic fixtures. No real brokerage observations were supplied to the new learning stage during development. The CVaR and local--stochastic-volatility sections below are a theoretical specification and literature synthesis; a CVaR hedge optimizer and an LSV calibration engine are not claimed as implemented.
\vspace{3mm}
\textbf{Use before submission.} Independently verify the mathematics, read the cited sources, replace synthetic diagnostics with a properly documented empirical study, record human contributions, and obtain both instructors' conditions for shared-course work and AI assistance.
\end{titlepage}

\begin{abstract}
This project investigates whether observable order-book information can improve the cost--risk tradeoff of derivative hedging, and whether any measured benefit survives changes in volatility dynamics. It separates three tasks that are often conflated: forecasting displayed liquidity under an observed information set; choosing a constrained hedge under a specified loss distribution; and valuing a derivative under a risk-neutral model. The implemented next stage constructs clock-time liquidity targets from stopped, integrity-checked displayed-depth exports. It uses only backward-as-of observations, records eligibility losses, splits whole receipt dates, and benchmarks persistence against regularized linear corrections. This report develops the corresponding measurement and estimation theory, derives a convex scenario-CVaR hedge formulation, and explains the role and limits of local--stochastic volatility. A critical review connects order-flow imbalance, order-book prediction, regularization, execution theory, conditional tail risk, particle calibration, and learned hedging. The proposed contribution is a reproducible comparison of information and decision rules, not a claim of a new pricing theorem or a profitable trading system. Synthetic validation supports software correctness and timing invariants; empirical forecasting and downstream hedging benefit remain open tests for the course project.
\end{abstract}
\tableofcontents
\newpage

\section{Research question, course mapping, and evidence}
\subsection{A decision problem rather than a collection of models}
The central question is:
\begin{quote}
Does order-book information improve held-out liquidity forecasts sufficiently to change constrained hedge decisions, and do those changes improve the hedging cost--risk frontier under explicitly specified volatility dynamics?
\end{quote}
The causal chain to investigate is measurement $\rightarrow$ prediction $\rightarrow$ scenario construction $\rightarrow$ decision $\rightarrow$ self-financing accounting. Each arrow requires a separate test. Better forecast accuracy does not logically imply better tail-risk control, and an option-calibrated process does not automatically supply real-world probabilities of loss.

The minimum experiment uses one underlying, one declared direct depth feed, one stock hedge instrument, and one hypothetical European option liability. A 30-second horizon and a fixed displayed trade quantity are initial research settings, not market constants. The code defaults to 100 shares and one-second sampling; pilot evidence must establish whether that quantity is observable before it is frozen for held-out evaluation. Passive queue-position fills, market impact, full option-chain calibration, and live order submission are outside the first scope.

\subsection{Source-grounded academic requirements}
The supplied AMS 518 syllabus requires each student to complete and present a numerical project. Its objectives include constrained optimization, transaction costs, financial systems with uncertain outcomes, and numerical case studies; listed topics include hedging, replication, quantile methods, and CVaR regression/estimation. Its project AI policy requires clear acknowledgement \citep[pp.~2--3]{ams518}. The proposed contribution is therefore an optimization case study, not only a forecast-error table.

The supplied AMS 520 syllabus requires an interactive Colab/Jupyter report, three presentations, and code/documentation in GitHub without data credentials. Phase I concerns literature, collection, and cleaning; Phase II concerns EDA and first empirical results; the stated milestones are Lecture 10, Lecture 22, and the final-exam day \citep[pp.~2--3]{ams520}. The proposed contribution is time-respecting liquidity prediction and its downstream decision value. The supplied materials do not establish permission for identical cross-course submissions, a specific AMS 520 AI policy, or the details of the separately referenced AMS 518 project requirements. Those gaps are not filled by inference.

\par\medskip
\begin{tabularx}{\textwidth}{@{}p{24mm}YY@{}}
\toprule
Layer & AMS 518 emphasis & AMS 520 emphasis\\\midrule
Shared input & Defined losses and liquidity constraints & Observed features, targets, and timing\\
Method & Convex scenario-CVaR optimization & Regularization and temporal generalization\\
Evidence & Feasible cost--risk comparisons & Held-out forecasts and information ablations\\
Submission & Distinct numerical case study & Interactive notebook and presentations\\\bottomrule
\end{tabularx}
\par\medskip

\subsection{What is established at this increment}
The base research platform at revision \texttt{8f35475b} records displayed-depth events and reconstructs them in C++ and Python. The new learning modules add feature/target alignment and persistence/ridge experiments. No new database schema, server route, order capability, or subscription purchase is required. The new pipeline reads exported files offline. The theoretical optimizer, conditional cost distribution, fitted nonlinear model, and LSV simulator/calibration remain subsequent work.

Throughout this report, \emph{source result} means a finding attributed to a cited work; \emph{derivation} means an argument shown here under stated assumptions; and \emph{design choice} means an implementation or proposed experiment specific to this project. Synthetic fixtures establish neither actual feed quality nor an out-of-sample market advantage.

\section{Critical literature review}
\subsection{Order-book state and order-flow information}
\citet{cont2014} study the relation between best-quote order-flow imbalance (OFI), depth, and short-interval price changes. Their parsimonious representation motivates an interpretable flow feature rather than an immediate dependence on a large neural network. The transfer to this project is limited: contemporaneous price impact is not the same target as future displayed trading cost. A variable constructed over the forecast interval would be unavailable when the forecast is made. Here OFI is accumulated only over a trailing interval ending at the decision time.

\citet{gould2016} examine queue imbalance for the direction of the next mid-price movement, using logistic models and distinguishing large- and small-tick behavior. This supports testing a simple imbalance baseline and checking instrument dependence. It does not establish that imbalance forecasts a 30-second spread, that five broker-provided rows reproduce an exchange queue, or that a directional classifier improves option hedging after costs. The project changes both the response variable and the observation mechanism; results must be re-established on its own data.

\citet{zhang2019} combine convolutional and recurrent components for order-book price prediction. Their work is a relevant nonlinear benchmark direction, not a reason to skip persistence, linear baselines, or data diagnostics. In this project a neural model would require sufficient independent sessions, stable feed conventions, and a predeclared validation protocol. A success on a published benchmark is not a performance guarantee for a limited direct-depth feed with local receipt timestamps.

\subsection{Regularization, validation, and dependence}
The shrinkage perspective of \citet{hoerl1970} motivates ridge regression when correlated features make unregularized coefficients unstable. The project uses an SVD-based solution of a specified mean-loss objective, preserving its scaling convention explicitly. Ridge is a controlled statistical baseline, not a structural model of exchange behavior. Its penalty is chosen on validation dates rather than on the final test interval.

Train-only preprocessing is necessary because transformations estimated on a test set expose information unavailable during fitting. The official scikit-learn guidance highlights this leakage mechanism \citep{sklearnpitfalls}. The implementation does not require scikit-learn: its scaler and ridge coefficients are explicit NumPy objects, so their dependence on training rows is directly testable.

\citet{diebold1995} develop forecast comparisons that allow serially correlated errors and general loss functions. \citet{politis1994} provide a resampling construction for weakly dependent stationary observations. These works motivate dependence-aware inference, but their assumptions cannot be replaced by a large raw event count. The current code reports per-date and pooled errors and paired per-date loss differences; it does not report IID significance tests. A block-bootstrap or long-run-variance analysis is future empirical work, with block choice and nonstationarity diagnostics documented rather than hidden.

\subsection{Liquidity costs and optimal execution}
\citet{almgren2001} separate transaction-cost effects from price-risk exposure over a liquidation horizon and study the tradeoff through an efficient frontier. The useful lesson here is architectural: execution conditions belong in the decision objective and constraints, not as an afterthought deducted from a frictionless hedge. However, an observed static book-crossing calculation is not an identified permanent or temporary impact model. This project initially excludes the feedback from its own trades to subsequent books. That small-agent assumption must be retained in any interpretation.

\subsection{Tail-risk optimization and learned hedging}
\citet{rockafellar2000} supply the auxiliary-threshold representation that makes scenario CVaR amenable to optimization. The project uses this representation in its proposed constrained hedge problem. It does not equate minimizing CVaR with minimizing VaR, optimizing the confidence level, or finding a globally optimal multistage strategy.

\citet{buehler2018} formulate learned hedging under trading frictions and convex risk objectives. This connects prediction to sequential decisions and motivates a later policy-learning comparison. It does not eliminate the need for correct cash accounting, nonanticipating observations, a defensible scenario distribution, or a feasible policy benchmark. Deep reinforcement learning is therefore an extension after the explicit constrained optimizer, not the dependency on which the initial submission rests.

\subsection{Local volatility and stochastic volatility}
\citet{dupire2010} describes the relation between European option prices across strike/maturity and local volatility. \citet{heston1993} provides a stochastic-volatility valuation framework with spot/variance correlation. The models address derivative valuation, not the empirical identification of order-book event intensities. Their role here is to generate controlled alternative derivative-risk environments while the hedge policy and liquidity assumptions are held inspectable.

The distinction between marginals and dynamics is fundamental. \citet{gyongy1986} establishes a mimicking result under stated hypotheses: suitable Markov coefficients can reproduce one-dimensional marginal distributions of an It\^o process. Equality of these marginals does not establish equality of transition laws, path-dependent prices, hedge errors, or the joint distribution with liquidity. Section~\ref{sec:lsv} derives the leverage identity while retaining these boundaries.

\subsection{LSV calibration is an inverse problem, not an input field}
\citet{cozma2019} develop a control-variate particle approach for a hybrid local--stochastic-volatility model with stochastic rates, building on the Guyon--Henry-Labord\`ere particle approach. This motivates a later numerical calibration experiment and attention to conditional-expectation error. The initial course target is narrower: deterministic rates and a controlled leverage function or a synthetic calibration target.

\citet{wyns2016} study an adjoint approach to matching a discretized local-volatility model. This provides an alternative to particle calibration and a reminder that forward and backward discretizations must be consistent. A finite-grid calibration identity is not automatically an error-free continuum calibration.

\citet{jourdain2020} address existence in a calibrated regime-switching setting. It should not be cited as an unrestricted existence theorem for every Heston-type leverage specification. \citet{bayer2024} study regularized conditional-expectation approximation through reproducing-kernel Hilbert spaces for singular LSV McKean--Vlasov problems. Together these works show why denominator regularization, conditioning diagnostics, and numerical stability belong in the LSV study; the review does not claim those algorithms have been reproduced in this repository.

\subsection{Research gap being tested, not novelty being asserted}
The selected literature justifies the ingredients but does not settle this project's question. Its contribution is a controlled link between a particular observable feed, future liquidity prediction, and constrained derivative hedging. The relevant ablation is not simply ``linear versus deep.'' It is whether top-of-book and additional recorded depth improve forecasts beyond current cost/price history, and whether that improvement remains useful once constraints and tail outcomes are evaluated.

This is a targeted critical review, not a systematic or exhaustive survey. No ``first ever'' claim is made. A negative result is scientifically useful when the data boundary, policy comparison, and uncertainty are reported correctly.

\begin{longtable}{@{}p{31mm}p{53mm}p{66mm}@{}}
\caption{Literature-to-experiment map. Each implication is a proposed test, not an imported empirical conclusion.}\label{tab:literature}\\
\toprule
Source family & Relevant mechanism & Project test / principal limitation\\\midrule
\endfirsthead
\toprule Source family & Relevant mechanism & Project test / principal limitation\\\midrule
\endhead
Cont et al.; Gould--Bonart & Flow and state imbalance contain different information & Compare trailing OFI with static imbalance; do not use forecast-interval flow\\
DeepLOB & Nonlinear spatial/temporal feature extraction & Deferred benchmark; small proprietary samples may not support the architecture\\
Hoerl--Kennard & Shrink correlated predictors & Tune ridge only on validation dates; coefficients are not causal estimates\\
Diebold--Mariano; Politis--Romano & Dependence-aware forecast comparison & Treat overlapping errors as dependent; per-date replication matters\\
Almgren--Chriss & Liquidity cost versus exposure risk & Do not relabel displayed-depth walking as a calibrated impact model\\
Rockafellar--Uryasev & Convex scenario tail-loss objective & Optimize holdings and threshold at fixed confidence, not the confidence itself\\
Deep Hedging & Risk-based policy learning with frictions & Compare against a transparent constrained policy and a reconciled ledger\\
Dupire; Heston; Gy\"ongy & Marginal calibration and alternative variance dynamics & Test model sensitivity; matching vanilla prices does not match hedge dynamics\\
Cozma et al.; Wyns--in 't Hout; Bayer et al. & Numerical conditional expectations and consistent calibration & Validate on a known synthetic target before fitting sparse market quotes\\
\bottomrule
\end{longtable}


\section{Measurement: what the recorder actually observes}
\subsection{Displayed depth is not an individual-order book}
Let the available reconstructed book be
\[
B_t=\{(b_{t,k},q^b_{t,k}),(a_{t,k},q^a_{t,k}): k=1,\ldots,K_t\},
\]
with descending bids, ascending asks, and nonnegative displayed sizes. The current adapter requests a bounded number of rows from one direct USD equity/ETF feed. Rows may refer to market makers; equal-price rows are aggregated when a top-price quantity is required. There is no individual-order identifier or verified passive queue position.

IBKR's documentation warns that its depth feed does not guarantee every quoted price and excludes odd lots. It also describes book clearing after depth-reset message 317 \citep{ibkrdepth}. Accordingly the code calls a structurally acceptable state \texttt{two\_sided\_unverified}. This name is substantive: it does not certify completeness, freshness of every row, or an executable fill. Requested rows are a bound, not proof that all those rows arrived.

\subsection{Observation filtration and clocks}
An archived event has a local sequence, a local monotonic receipt time $\tau_i$, a wall-clock receipt coordinate $u_i$, and an operation/side/row payload. Define
\[
\G_t=\sigma\{E_i:\tau_i\leq t\}.
\]
The reconstruction at time $t$ must be a function of $\G_t$. This is an \emph{observed-feed filtration}, not the exchange's information filtration. Feed latency, aggregation, missing upstream events, and process scheduling can separate the two.

A regular monotonic grid $t_j=t_0+j\delta$ avoids giving busy periods more regression weight merely because more events arrived. The last eligible index is
\begin{equation}
 i(t)=\max\{i:\tau_i\leq t\}.\label{eq:asof}
\end{equation}
All events with a tied receipt time are applied in archived sequence before sampling that instant. A forward or nearest-time join is forbidden. Wall time is used for date grouping and time-of-day features; duration is measured with the monotonic clock. A backward clock is rejected. A material inconsistency between clock increments breaks the continuity segment.

\begin{proposition}[Prefix nonanticipation]
Suppose reconstruction is deterministic, grid evaluation uses \eqref{eq:asof}, and features depend only on reconstructed observations in $[t-L,t]$. Altering events whose receipt times exceed $t$ cannot alter those feature values at $t$.
\end{proposition}
\begin{proof}
Every selected event index is at most $i(t)$. Deterministic reconstruction of that prefix is unchanged by an alteration outside it. Each trailing feature is a deterministic function of unchanged prefix states. No statement is made about whether a future-dependent label remains eligible after the alteration.
\end{proof}

The last sentence matters. A model feature can be causal while the \emph{evaluation sample} is selected using future data availability. Those are different issues, addressed below.

\subsection{Continuity and age rules}
A reset, gap, terminal marker, or invalid event state creates a new segment or invalidates the state. A lookback or prediction horizon cannot cross such a break, even if the next sampled book appears plausible. No interpolation recreates deleted or missing rows.

The pipeline also requires both sides to have received an update within an explicit bound $A$ seconds at sampled times. A new ask update cannot refresh the bid-side clock. This is only a \emph{side-update age} rule: a deep bid update can make the side's clock recent while a different bid row is old. Row-level age, exchange-time synchronization, and true quote lifetime remain unverified. The default $A=5$ seconds is a conservative experiment setting that must be reported and sensitivity-tested.

\section{Liquidity targets and observable features}
\subsection{Spread and displayed crossing cost}
For an uncrossed, unlocked two-sided state, define
\[
 m_t=\frac{a_{t,1}+b_{t,1}}{2},\qquad s_t=a_{t,1}-b_{t,1}.
\]
For a fixed positive quantity $Q$, let $x^+_{t,k}(Q)$ be the amount consumed at ask row $k$ when the rows are walked in order, with $0\leq x^+_{t,k}\leq q^a_{t,k}$ and $\sum_kx^+_{t,k}=Q$. Define $x^-_{t,k}$ similarly for a sale. The dollar proxies are
\begin{align}
 c_t^+(Q)&=\sum_k a_{t,k}x^+_{t,k}(Q)-Qm_t,\\
 c_t^-(Q)&=Qm_t-\sum_k b_{t,k}x^-_{t,k}(Q).
\end{align}
They exist only when the recorded side has sufficient size. The normalized quantities are
\begin{equation}
 z_t^{\pm}(Q)=10^4\frac{c_t^{\pm}(Q)}{Qm_t},\qquad
 z_t^{s}=10^4\frac{s_t}{m_t}.\label{eq:costbps}
\end{equation}
The code supports a chosen one of \texttt{buy\_cost\_bps}, \texttt{sell\_cost\_bps}, or \texttt{spread\_bps}. Its default primary response is
\[
 Y_t=z_{t+H}^{+}(Q),\qquad H=30\text{ seconds}.
\]
The midpoint in the label denominator is the midpoint at the label time. It is not a future input to the predictor; it is part of the definition of the future outcome.

These proxies exclude delay to order arrival, disappearing displayed size, hidden liquidity, price impact, broker fees, and actual execution mechanics. They measure future \emph{displayed conditions}. They are not transaction-cost observations from executed trades. Nor is forecasting $z_t$ from a contemporaneous book a genuine forecast; persistence uses that directly observable quantity as its benchmark.

\subsection{Static imbalance and trailing flow}
Aggregate all rows at the best bid and ask to obtain $q_t^{b,0}$ and $q_t^{a,0}$. Top imbalance is
\[
 I_t^0=\frac{q_t^{b,0}-q_t^{a,0}}{q_t^{b,0}+q_t^{a,0}}.
\]
Total recorded-depth imbalance replaces these top quantities with sums over all observed rows. ``Total'' here always means total \emph{recorded} depth.

For consecutive structurally eligible events with best prices $b_i,a_i$ and top quantities $q_i^b,q_i^a$, the implemented OFI increment follows the best-quote construction of \citet{cont2014}:
\begin{align}
 e_i={}&\ind_{\{b_i\geq b_{i-1}\}}q_i^b-\ind_{\{b_i\leq b_{i-1}\}}q_{i-1}^b\nonumber\\
 &-\ind_{\{a_i\leq a_{i-1}\}}q_i^a+\ind_{\{a_i\geq a_{i-1}\}}q_{i-1}^a.
\end{align}
It is accumulated on $(t-L,t]$ and divided by the current total best-price quantity. It is not signed trade volume: quote changes need not identify whether an execution or cancellation occurred. In this implementation the flow accumulation restarts after invalid states.

\subsection{Predeclared information sets}
Every model uses the same eligible examples. The nested groups are:
\par\medskip
\begin{tabularx}{\textwidth}{@{}p{25mm}Y@{}}\toprule
Group & Inputs available at decision time\\\midrule
History & Current chosen target, trailing midpoint log return, sample standard deviation of trailing grid log returns, UTC time-of-day sine/cosine\\
Top & History plus spread, best-price imbalance, log best bid/ask quantities, scaled trailing OFI\\
Depth & Top plus recorded-depth imbalance and log total bid/ask quantities\\\bottomrule
\end{tabularx}
\par\medskip
The history group is therefore \emph{not} a strictly price-only baseline: it includes the observed current liquidity target. This is intentional, because incremental prediction should beat an informative persistence benchmark. Size transformations use $\log(1+q)$. Trailing return standard deviation is unannualized and is not passed into a pricing model.

\subsection{Future availability creates a selected population}
Let $A_t$ denote eligibility: valid and sufficiently recent states throughout the feature and horizon windows, the same reconstruction segment and UTC date, and enough displayed size at the endpoints for the chosen target. The first regression estimates
\[
 \E[Y_t\mid X_t=x,A_t=1],
\]
not automatically $\E[Y_t\mid X_t=x]$. Future capacity or outages can depend on stress. Dropping precisely those intervals can make the measured sample easier and conceal economically important tail cases.

The output therefore records disjoint exclusion counts: warmup/right boundary, invalid or stale features, invalid or stale horizon, cross-date windows, insufficient current depth, insufficient future depth, and eligible rows. Counts reconcile to the clock grid. This is transparency, not a statistical correction. A later availability model and a conservative policy for unpriced depth are needed before claiming an unconditional decision benefit. No missing label is changed to zero, and no prediction can justify trading beyond observed capacity.

\section{Regularized forecasting and temporal validation}
\subsection{A correction to persistence}
Write $y_i$ for the future target and $p_i$ for its contemporaneous value. Persistence is $\widehat y_i=p_i$. Ridge models the residual $d_i=y_i-p_i$:
\begin{equation}
 \min_{b,\beta}\frac1n\sum_{i\in\mathcal T}(d_i-b-z_i^\top\beta)^2
              +\lambda\norm{\beta}_2^2,\quad\lambda>0.\label{eq:ridge}
\end{equation}
Only training rows $\mathcal T$ determine feature means and population-standard-deviation scales. Constant or numerically near-constant scales are replaced by one; no future row determines that decision. The intercept is unpenalized. The positive penalty controls unstable directions, in the spirit of \citet{hoerl1970}.

If $Z=U\operatorname{diag}(s_j)V^\top$ is the training SVD and $d_c=d-\bar d\mathbf1$, then
\begin{equation}
 \widehat\beta=V\operatorname{diag}\left(\frac{s_j}{s_j^2+n\lambda}\right)U^\top d_c,
 \qquad \widehat b=\bar d.\label{eq:svd}
\end{equation}
This follows by differentiating \eqref{eq:ridge} and resolving the positive-definite quadratic in singular-vector coordinates. The implementation uses \eqref{eq:svd}, not an explicit inverse of $Z^\top Z$. A library whose objective uses a \emph{sum} rather than a mean would require penalty $n\lambda$ for the same fit.

Predictions are $p_i+\widehat b+z_i^\top\widehat\beta$. Negative predictions are counted, not silently clipped. They are evidence of a mismatch between an unconstrained conditional-mean baseline and a nonnegative response. Any later clipping or constrained regression must be predeclared and evaluated consistently, not chosen after viewing test errors.

\subsection{Whole-date partitions and information intervals}
All captures on one UTC receipt date belong to a single partition. Training dates precede validation dates, which precede test dates. For each eligible example define its complete information/label interval
\[
 J_t=[t-L,t+H].
\]
The builder excludes windows crossing a receipt date. The splitter additionally checks that the last earlier-partition label precedes the first later-partition feature-window start. No random row split or row-count gap substitutes for this time check.

Real-data manifests must explicitly declare every eligible date in exactly one of the three partitions. Synthetic mode uses a deterministic date allocation only for integration tests. Three dates are a software minimum, not adequate evidence of generalization. Dates with no usable labels require investigation; the code refuses an incompatible explicit partition instead of inventing observations.

Within each predeclared information set, $\lambda$ is selected on equal-date validation MSE from a fixed candidate list. The final feature-family selection also uses validation only; all declared families are shown on the test set as ablations. The initial implementation retains the training-only fits rather than refitting on validation. Repeatedly changing the target, quantity, filter, or model after reading test scores would invalidate the holdout even though the code's initial split was correct.

\begin{proposition}[Test-set noninterference in fitting]
With a fixed partition and candidate set, the fitted scalers and coefficients in \eqref{eq:ridge}, and a selection rule depending only on validation losses, are invariant to changes in test features or labels.
\end{proposition}
\begin{proof}
The scaler and SVD depend only on training arrays. Validation predictions depend on those fixed estimates and validation features; validation losses use validation labels. None of those functions has the test arrays as an argument. The test arrays enter only the final scoring operation. This is tested by perturbing test labels and feature magnitudes.
\end{proof}

\subsection{Metrics and what they do not prove}
The code reports pooled MAE, RMSE, bias, negative-prediction fraction, and per-date errors. Hyperparameter selection uses
\[
 \operatorname{MacroMSE}=\frac1D\sum_{d=1}^{D}\frac1{n_d}\sum_{i\in d}(\widehat y_i-y_i)^2.
\]
This gives each date equal evaluation weight even when retention or collection duration differs. Model fitting remains row-weighted, an explicit distinction. For each test date it also reports the paired squared-error difference relative to persistence on the same examples.

With $H=30$ and $\delta=1$, adjacent forecasts share much of the same horizon and trailing history. The observations are not IID. Date aggregation reduces one source of pseudo-replication but does not establish independence across days. Dependence-aware tests and resampling require their own assumptions \citep{diebold1995,politis1994}. No $p$-value or claimed confidence interval for improvement is emitted by this baseline implementation.


\section{Derivative risk: local and stochastic volatility}\label{sec:lsv}
\subsection{Separate the pricing and scenario measures}
Let $\Q$ be a selected pricing measure and $\Pp$ the distribution used for physical-risk or controlled scenario evaluation. A risk-neutral European value is
\[
 V(t,s,v)=\E^{\Q}\!\left[e^{-\int_t^T r_u\,du}g(S_T)\mid S_t=s,v_t=v\right].
\]
A risk calculation such as $\CVaR_\alpha^{\Pp}(L)$ requires an explicitly specified $\Pp$ distribution. A fitted option surface does not identify the variance risk premium or all physical drifts. Using identical parameters under both measures is permissible as a \emph{declared controlled experiment}, not as an empirical conclusion.

\subsection{Local volatility and the option surface}
For deterministic rates/yields and a sufficiently smooth, arbitrage-consistent call surface $C(K,T)$, the local-variance expression is
\begin{equation}
 \sigma_{\mathrm{loc}}^2(T,K)=
 \frac{C_T+(r_T-q_T)K C_K+q_T C}{\tfrac12 K^2 C_{KK}}.\label{eq:dupire}
\end{equation}
The forward relation and local interpretation are reviewed by \citet{dupire2010}. The denominator is related to a risk-neutral density; noisy second strike derivatives are not harmless input noise. A candidate empirical surface needs convention checks, monotonicity/convexity controls, maturity consistency, and a disclosed smoothing procedure. This project does not claim an arbitrage-free market surface has been collected or fitted.

\subsection{LSV dynamics and the leverage relation}
Consider the representative model
\begin{align}
 \frac{dS_t}{S_t}&=(r_t-q_t)dt+\ell(t,S_t)\sqrt{v_t}\,dW_t^{S,\Q},\\
 dv_t&=\kappa(\theta-v_t)dt+\xi\sqrt{v_t}\,dW_t^{v,\Q},\qquad
 d\langle W^{S,\Q},W^{v,\Q}\rangle_t=\rho\,dt.
\end{align}
The stochastic-variance component is Heston-type \citep{heston1993}; the leverage function $\ell$ supplies a local component. Conditional on $S_t=s$, instantaneous stock variance is
\[
 s^2\ell^2(t,s)\E^{\Q}[v_t\mid S_t=s].
\]
Matching this conditional variance to the projected local model yields the calibration identity
\begin{equation}
 \ell^2(t,s)\E^{\Q}[v_t\mid S_t=s]=\sigma_{\mathrm{loc}}^2(t,s).\label{eq:leverage}
\end{equation}
The projection interpretation requires suitable hypotheses; it is not an unconditional existence statement for arbitrary coefficients \citep{gyongy1986,jourdain2020}.

Equation~\eqref{eq:leverage} is nonlinear in the model law: the conditional expectation depends on trajectories generated using $\ell$ itself. One kernel approximation would take
\[
 \widehat m_\varepsilon(t,s)=
 \frac{\sum_{j=1}^N v_t^{(j)}K_\varepsilon(S_t^{(j)}-s)}
      {\sum_{j=1}^N K_\varepsilon(S_t^{(j)}-s)}.
\]
Sparse tails can make this estimate unstable. Particle variance reduction, consistent PDE methods, and regularized conditional-expectation approximations address related numerical problems \citep{cozma2019,wyns2016,bayer2024}. Bandwidths, denominator floors, boundary handling, and variance-discretization choices must be reported as numerical assumptions, with calibration error measured afterward. A positive floor is not proof of exact calibration.

\subsection{What a single stock hedge cannot remove}
For a sufficiently regular $V$, its diffusion exposure under the displayed model contains
\[
 V_s\,S\ell\sqrt v\,dW^S+V_v\,\xi\sqrt v\,dW^v.
\]
Writing $dW^v=\rho dW^S+\sqrt{1-\rho^2}\,dW^\perp$ shows that a stock-only minimum local-variance projection, under the specified instantaneous covariance, is
\begin{equation}
 h^{\mathrm{proj}}=V_s+\frac{\rho\xi}{S\ell}V_v,\label{eq:hedgeprojection}
\end{equation}
when $S\ell\sqrt v\ne0$. It follows by minimizing the variance of the option diffusion minus $h\,dS$. The unspanned component remains
$V_v\xi\sqrt{v(1-\rho^2)}dW^\perp$.

This derivation is an \emph{instantaneous quadratic-risk} result, not the constrained CVaR solution below. A chosen pricing delta $V_s$ and a minimum-variance projected hedge need not agree. Under different physical covariance assumptions the relevant projection also changes. An oracle experiment may use the simulated $v_t$; a practically observable policy may not silently receive that latent state. Estimated variance and true simulated variance must be separate information sets.

Matching European price marginals via \eqref{eq:leverage} cannot guarantee the same hedge-error law. Nor does it identify any joint relationship between variance, spread, and depth. A controlled LSV study should initially hold a declared liquidity model fixed, then stress dependence rather than attaching unrelated historical books to arbitrary simulated prices and calling the result empirical.

\section{Self-financing hedging and the CVaR decision layer}
\subsection{Account for trades before optimizing them}
For one hypothetical short option, let $W_t$ be stock-plus-cash hedge wealth just before a trade, $h_{\mathrm{old}}$ the previous holding, $h$ the new holding, and $u=h-h_{\mathrm{old}}$. Assume $q=0$, fractional shares, symmetric funding at a deterministic rate, and exogenous scenarios. These simplify the first case study; they are not assertions about broker financing.

After paying current trading cost $C_t(u)$, the cash account is
\[
 B_t^+=W_t-C_t(u)-hS_t.
\]
For $R=e^{r\Delta}$ and scenario $j$ with stock $S'_j$, marked option liability $V'_j$, and future stock-liquidation cost $C'_j(-h)$, the horizon loss is
\begin{align}
 L_j(h)&=V'_j-\{R B_t^+ +hS'_j-C'_j(-h)\}\nonumber\\
       &=\underbrace{V'_j-RW_t}_{b_j}
          -h\underbrace{(S'_j-RS_t)}_{g_j}
          +R C_t(h-h_{\mathrm{old}})+C'_j(-h).\label{eq:loss}
\end{align}
For receding-horizon use, liquidation in this expression is a horizon valuation convention. The realized policy ledger charges only trades actually executed; a hypothetical liquidation cost must not be debited at every replanning step unless liquidation occurs. A terminal horizon uses the payoff in place of $V'_j$. Interior liability marks and terminal settlement must not both be charged as cash payments. No hidden cash injection makes an infeasible hedge appear funded.

The future forecast has a meaningful role in $C'_j$, not in replacing the observed cost $C_t$ of a trade now. A 30-second mean-cost prediction is not enough to determine the joint scenario law of $(S'_j,V'_j,C'_j)$. The first learning stage establishes only a point-prediction interface; conditional residual scenarios and dependence calibration remain to be implemented.

\subsection{Convex costs from displayed rows}
For a static ordered book and zero fixed fees, the displayed cost of buying quantity $u\geq0$ is the integral of nondecreasing ask-minus-mid marginal prices. Selling $u<0$ gives an analogous signed cost. On the available-size domain the combined function is convex and piecewise linear: its slope increases from negative bid-side slopes to positive ask-side slopes.

\begin{proposition}[Convex scenario hedge problem]
Suppose current and horizon costs are convex functions of the signed stock trades, scenario values and probabilities are exogenous, and the feasible holding set is convex. Then each $L_j(h)$ in \eqref{eq:loss} is convex. The scenario-CVaR minimization in $(h,\eta)$ is jointly convex.
\end{proposition}
\begin{proof}
The first two terms in \eqref{eq:loss} are affine. Positive scaling preserves convexity of $C_t$, and composition with an affine map preserves convexity of both costs. For each $j$, $\max\{L_j(h)-\eta,0\}$ is the maximum of two convex functions. A nonnegative weighted sum plus the affine function $\eta$ is convex.
\end{proof}
This is an application of standard convex composition rules \citep{boyd2004}. Nonlinear option values across scenarios do not make dependence on the \emph{decision} $h$ nonconvex when those scenario liabilities are fixed. Conversely, a fixed fee per nonzero trade, integer lots, decision-dependent market impact, or an unrestricted learned cost curve can invalidate the continuous convex formulation.

\subsection{CVaR and the auxiliary threshold}
For integrable loss $L$ and fixed $0\leq\alpha<1$, use
\begin{equation}
 \CVaR_\alpha(L)=\inf_{\eta\in\mathbb R}
 \left\{\eta+\frac{1}{1-\alpha}\E[(L-\eta)^+]\right\}.\label{eq:cvar}
\end{equation}
For $0<\alpha<1$, an integrable real-valued loss has finite quantile minimizers. At $\alpha=0$, the infimum equals $\E[L]$ but need not be attained by a finite threshold when losses are unbounded below; a finite empirical sample does attain it. The representation underlies scenario optimization \citep{rockafellar2000}. The confidence level $\alpha$ is chosen; $\eta$ is optimized and need not be unique. Atoms require tail-mass allocation, so $\E[L\mid L\geq\VaR_\alpha]$ is not a generally interchangeable definition.

For fixed $h$, the subgradient interval in $\eta$ is
\[
 \left[\frac{\Pp(L(h)<\eta)-\alpha}{1-\alpha},
       \frac{\Pp(L(h)\leq\eta)-\alpha}{1-\alpha}\right].
\]
Thus a minimizer satisfies $\Pp(L(h)<\eta)\leq\alpha\leq\Pp(L(h)\leq\eta)$. For example, loss zero with probability $.95$ and loss ten with probability $.05$ has lower $.95$-quantile zero and CVaR ten; every $\eta\in[0,10]$ minimizes the auxiliary expression. A unique threshold should not be inferred from an optimizer's arbitrary selection within this interval.

With equal-probability scenarios, a proposed AMS 518 problem is
\begin{align}
 \min_{h,\eta,z}\quad&\eta+\frac{1}{N(1-\alpha)}\sum_{j=1}^N z_j\label{eq:lp}\\
 \text{s.t.}\quad& z_j\geq L_j(h)-\eta,\quad z_j\geq0,\nonumber\\
 &h_{\min}\leq h\leq h_{\max},\quad |h-h_{\mathrm{old}}|\leq Q_{\max},\nonumber\\
 &C_t(h-h_{\mathrm{old}})\leq B.\nonumber
\end{align}
Piecewise-linear costs can be represented with epigraph variables to obtain an LP. $Q_{\max}$ should not exceed the displayed-size domain without an additional cost model. Holdings must also remain in the domain of each modeled future liquidation cost, or an additional conservative cost/capacity treatment must be specified. Cardinality or nonzero-trade fees require a separately stated mixed-integer or nonconvex extension.

\subsection{Why predicting mean liquidity is not predicting tail risk}
Consider a cost variable that is always one, and another that is zero with probability $.9$ and ten with probability $.1$. Both have mean one, but their $.95$-CVaRs are one and ten. No mean predictor can distinguish their tail implications from that mean alone. Similarly, independent sampling of price and liquidity shocks can miss joint stress precisely where a hedge needs liquidity.

Future work must therefore specify a conditional residual law or scenario generator, include adverse dependence sensitivity, and measure calibration on later dates. A learned cost model should preserve the intended monotonic/convex quantity structure or supply scenario coefficients to a known convex family. Calling a repeated one-step solution ``optimal terminal CVaR hedging'' would also overstate the result: the proposed policy is receding-horizon unless a multistage formulation and time-consistency treatment are provided.


\section{Experimental design and falsifiable comparisons}
\subsection{Separate empirical and controlled tracks}
The \textbf{empirical track} records later direct-depth sessions, freezes compatible exports, learns future displayed-cost targets, and compares persistence with predeclared information sets. The \textbf{controlled track} uses declared price/variance and liquidity dynamics to test hedge policies with the existing self-financing ledger. An empirical regression coefficient does not by itself validate the controlled joint simulator.

The initial liquidity experiment holds target, quantity, time step, horizon, age bound, and split fixed. Its main ablation compares history, top, and depth on common eligible examples. A sensitivity analysis can later vary $Q$, $H$, and $A$, with an explicit separation between exploratory and confirmatory datasets. Recording many adjacent events is not a substitute for observing diverse sessions.

For hedging, compare periodic delta, a declared no-trade-band benchmark, baseline-liquidity CVaR, and book-informed CVaR. Hold liability, initial funding, constraints, and evaluation scenarios fixed when comparing decision rules. A different hedge volatility changes initial model premium in the current synthetic lab; this funding effect must not be misreported as policy skill.

\subsection{Primary metrics and failure modes}
Forecast metrics should include per-date MAE/MSE, paired improvements, negative predictions, retention, and performance in wide-spread or shallow-depth states defined using training/pilot rules. The terminal hedge study should include loss CVaR, mean loss, turnover, financed costs, binding constraints, and ledger residuals. Neither a histogram quantile nor a finite-sample CVaR estimate is a certified population tail bound.

The following failure modes are testable before empirical claims:
\par\medskip
\begin{tabularx}{\textwidth}{@{}p{37mm}Y@{}}\toprule
Failure & Required diagnostic or counterexample\\\midrule
Future event leaks into a feature & Alter future updates and verify the earlier feature prefix is unchanged\\
Test set changes the fit & Perturb test labels/features and verify coefficients and validation selections are unchanged\\
Busy dates dominate evaluation & Report equal-date metrics beside pooled metrics\\
Missing depth conceals stress & Reconcile exclusion counts and analyze availability by observed state\\
Reset crosses a label window & Reject the entire affected window, including between-grid invalidation\\
Predictive improvement lacks decision value & Evaluate identical feasible policies/scenarios with and without the information\\
Model calibration mistaken for physical forecast & State $\Q$ and scenario-law assumptions separately\\
Displayed cost mistaken for execution & Preserve the proxy label; do not infer queue position, impact, or fills\\\bottomrule
\end{tabularx}
\par\medskip

\subsection{What the synthetic run demonstrates}
The public fixture generates nine short synthetic dates with structured, mean-reverting latent liquidity conditions and small midpoint changes. Those dynamics are deliberately constructed, not estimated from real market data. Under the default 30-second lookback and horizon, 2,691 examples survive the rules in the tested local run. The deterministic partition uses five training dates, two validation dates, and two test dates. These are integration-test dimensions, not a sample-size recommendation.

The notebook prints and plots forecast diagnostics and verifies exact repeatability within the same environment. It does not display those numbers as empirical evidence. Software tests compare the SVD solution with an independently solved normal equation, verify timing and sample-count identities, reject incompatible or duplicate feeds, and check private no-overwrite outputs. Platform/library floating-point differences may change least-significant numerical digits; source and environment fingerprints are retained.

\section{Repository deliverables and reproducibility}
\subsection{Files added by this stage}
\par\medskip
\begin{tabularx}{\textwidth}{@{}p{63mm}Y@{}}\toprule
Component & Responsibility\\\midrule
\texttt{liquidity\_dataset.py} & Receipt-time reconstruction, clock sampling, features, labels, eligibility audit\\
\texttt{liquidity\_baselines.py} & Date split, train-only scaling, SVD ridge, validation selection, per-date scoring\\
\texttt{tools/liquidity\_study.py} & Explicit synthetic/private-manifest input and protected result export\\
\texttt{02\_liquidity\_prediction.ipynb} & Executable Phase II foundation with diagnostics and interpretation\\
\texttt{tests/liquidity/} & Counterexamples to timing leakage, estimator checks, notebook execution\\
\texttt{reports/theory\_literature.tex} & This reproducible theoretical and literature document\\\bottomrule
\end{tabularx}
\par\medskip
The new learner reads stopped exports and does not need a running server, IBKR credentials, or a second database. It does not modify the existing market archive. Its results are explicit private JSON files, not a silently added experiment type in the server catalog. No database migration is introduced by this stage.

A real-data manifest lists absolute local export paths, configuration, and all three date partitions. The command refuses mixed feed conventions, duplicate fingerprints, overlapping capture sessions, insufficient dates, and incompatible partitions. Outputs are owner-readable files published without overwrite outside detected Git checkouts. Market-data redistribution rights still require human review. The supplied syllabi and their meeting credentials are not copied into the repository.

\subsection{Run the synthetic demonstration}
From the existing checkout after obtaining the learning branch:
\begin{verbatim}
python3 -m pip install -r \
  research/liquidity_aware_hedging/requirements.txt
python3 tools/liquidity_study.py --synthetic \
  --output "$HOME/.local/share/derivative-lab/learning/demo.json"
python3 -m jupyterlab \
  research/liquidity_aware_hedging/notebooks/02_liquidity_prediction.ipynb
\end{verbatim}
A pre-existing output path is an error, not permission to overwrite an old run. To run real data, provide \texttt{--manifest /absolute/path/manifest.json} instead of \texttt{--synthetic}. For a private notebook copy, set \texttt{DTS\_REPO\_ROOT} and \texttt{DTS\_LIQUIDITY\_MANIFEST} before starting its kernel. Real-data mode never silently falls back to synthetic observations.

Compile this document with \texttt{latexmk -pdf -interaction=nonstopmode -halt-on-error theory\_literature.tex} from its report directory. The adjacent bibliography supplies the references. The PDF is a working-paper artifact, not a substitute for the AMS 520 notebook or either required presentation.

\section{Conclusion and next validation gate}
The immediate deliverable is a reproducible baseline for future displayed liquidity, not an end-to-end trading strategy. Its most important properties are explicit observation timing, common-sample comparisons, training-only estimation, predeclared date partitions, visible missingness, and reproducible output. These make a negative result interpretable and a positive result worth further testing.

The next empirical gate is multiple compatible recorded dates with a frozen manifest. The next optimization gate is a constrained CVaR solver with independently checked cash/stock losses and exogenous convex costs. The next volatility gate is a controlled stochastic-volatility/LSV experiment with a validated pricing reference and clear $\Pp/\Q$ assumptions. Only after those gates should the project claim that a liquidity forecast improves the cost--risk frontier, or extend to nonlinear learning and sequential policy optimization.

\section*{AI assistance, authorship, and source limitations}
\addcontentsline{toc}{section}{AI assistance, authorship, and source limitations}
This draft and the associated implementation were prepared with substantial AI assistance: research framing, source discovery and synthesis, derivations, Python code, test design, notebook drafting, and LaTeX production. The base repository also contains AI-assisted work. Automated checks are distinct from the student's independent verification and ability to defend the submitted analysis. Michael Perez should review and revise the report and maintain a dated human-contribution log before submission. No instructor approval, teammate contribution, real-data result, or literature novelty is invented.

The syllabus mapping uses the two user-supplied Fall 2026 PDFs only for stated course requirements. The technical literature is externally researched and cited separately. This review reads selected primary papers and official documentation; it is not an exhaustive review of every recent LSV or microstructure method. In particular, theoretical assumptions, data licensing, physical-risk calibration, and actual depth entitlements remain material unresolved requirements for empirical deployment.


\clearpage
\bibliographystyle{plainnat}
\bibliography{literature}
\end{document}
